\chapter{Introduction}

\section{Background of the Study}
During a crisis, the core logistical challenge encountered by humanitarian efforts globally is defined as the Resource Allocation Problem (RAP) \cite{starr2014, knott1987}. The primary objective of the RAP is the equitable allocation of scarce life-sustaining resources, including food, water, and medical supplies, under conditions of intense time pressure and environmental uncertainty \cite{noyan2018}. Because disaster typologies vary widely—from sudden-onset events like earthquakes to slow-onset crises like droughts—there is a critical need for generalized global models capable of adapting to diverse allocation constraints.

Unlike commercial supply chains that are optimized for profit maximization and cost reduction, humanitarian operations are strictly driven by the need to minimize human suffering and maximize survival rates \cite{gutjahr2018}. This theoretical boundary dictates that systems designed for disaster relief must prioritize equity and impact over sheer efficiency or financial profit \cite{huang2012}. The fundamental variables involved in this complex setup include demand points located in calamity areas, supply points acting as centralized hubs, and constraints regarding stock availability and proportional assignment \cite{lin2011, cui2010}.

To properly evaluate the fairness and efficiency of these allocation efforts, resource allocation criteria are broadly classified into two categories: ordinal and cardinal \cite{vitoriano2011, ferrer2018}. Ordinal criteria evaluate allocations based on relative preferences and include metrics such as Pareto-efficiency, envy-freeness, and max-min fair share \cite{gutjahr2018, huang2012}. Cardinal criteria, on the other hand, rely on quantifiable utility values and encompass approaches such as max utilitarian, pure egalitarian, max leximin, and max Nash \cite{vitoriano2011}. While all of these criteria offer distinct theoretical frameworks for allocation, humanitarian operations require a specific focus on criteria that prioritize the most vulnerable populations. 

\section{Statement of the Problem}
Efficient crisis response often hinges on the underlying decision-making architecture \cite{balcik2010}. Decentralized efforts, or those lacking clear communication, often suffer from the ``Bullwhip Effect,'' a phenomenon where a lack of communication leads to the over- or underestimation of needs, resulting in either wasted resources or critical shortages \cite{balcik2010}. Furthermore, while existing computational methodologies often rely on fast heuristics to solve real-time crises, these approaches frequently neglect the intricate ethical needs for equity \cite{wang2018}. Exact mathematical approaches, such as Mixed-Integer Linear Programming (MILP), offer the analytical rigor required to solve these ethical dilemmas, yet researchers face significant challenges in mathematically defining and operationalizing complex equity metrics \cite{ni2015}.

Crucially, there is a notable gap in existing literature regarding generalized models that can dynamically resolve this bottleneck through exact mathematical means. Very few existing systems successfully formulate a multi-stage exact MILP that simultaneously integrates a combination of distinct fairness criteria such as max-min equity, deprivation costs, or proportional fairness to build an accurate and rigorously equitable allocation system.

In the local context of Davao City, the necessity for such a dynamic and centralized allocation framework is highly pronounced. The city is persistently vulnerable to severe natural hazards, particularly catastrophic flash flooding driven by the overflowing of major catchments like the Davao and Talomo River Basins, alongside frequent seismic activities \cite{cayamanda2021}. During these localized crises, the region's decentralized barangay-level relief efforts and local Disaster Risk Reduction and Management Operations Centers (DRRMOCs) frequently encounter severe logistical and coordination bottlenecks, resulting in inefficient resource augmentation and data interpretation challenges \cite{samson2024}. This fragmented communication inevitably triggers a localized Bullwhip Effect, where highly accessible urban centers receive redundant aid while isolated, severely affected communities experience critical shortages. Implementing a two-stage exact MILP model tailored to centralize and mathematically guarantee equitable distribution across multiple fairness metrics is therefore essential to enhancing Davao City's localized disaster resilience.

This study investigates the development of a centralized decision framework that integrates exact optimization techniques with multi-criteria fairness constraints to address these operational and ethical gaps. Specifically, the study aims to answer the following questions:
\begin{enumerate}
\item What are the specific constraints and historical data parameters that influence resource allocation during a generalized global disaster?
\item How can a centralized information system mitigate the Bullwhip Effect in humanitarian aid assignment?
\item How can distinct fairness criteria be mathematically combined and quantified within an exact optimization model to prioritize equity?
\item How can a two-stage Mixed-Integer Linear Programming (MILP) framework be designed to formulate a mathematically tractable resource allocation system optimized for deprivation cost minimization under multi-criteria equity post-evaluations?
\end{enumerate}

\section{Objectives of the Study}
The general objective of this study is to develop a generalized, automated resource allocation system for humanitarian aid utilizing a two-stage exact Mixed-Integer Linear Programming (MILP) framework that prioritizes multi-criteria equity.

The specific objectives that the research aspires to achieve are:
\begin{itemize}
\item To define the parameters of demand points in calamity areas and the operational constraints of centralized supply hubs across varying disaster typologies.
\item To design a Centralized Decision Framework based on a ``Hub-and-Spoke'' model to maintain a single source of truth for resource assignment and mitigate the Bullwhip Effect.
\item To apply Deprivation Cost Functions to mathematically represent and penalize how human suffering grows over time within the optimization model.
\item To implement Max-Min Fairness strategies as strict mathematical constraints aimed at prioritizing and assigning the most benefit to the areas that are most in need.
\item To formulate and solve a two-stage exact MILP model that seamlessly integrates these distinct fairness criteria to compute rigorous, accurate, and optimal resource allocation decisions.
\end{itemize}

\section{Significance of the Study}
The output of this study provides a critical operational advantage for humanitarian disaster relief efforts \cite{starr2014}. By utilizing modern information visibility tools within a centralized system, the proposed framework maintains a real-time ``single source of truth.'' This architectural choice is essential for preventing local hoarding and ensuring that the assignment of aid is calculated based on actual aggregated data rather than panic-driven requests \cite{sahin2023, wagenaar2017}. By integrating multiple fairness criteria within an exact two-stage MILP model, this research argues for a mathematically rigorous system design that systematically minimizes population-weighted deprivation costs, ensuring objective, hazard-weighted resource allocation.

On a localized level, this framework provides a targeted solution to the systemic logistical challenges faced by Davao City's Disaster Risk Reduction and Management Operations Centers (DRRMOCs) \cite{samson2024}. During catastrophic events, such as severe flash flooding in the Davao and Talomo River Basins, the city's reliance on decentralized, barangay-level relief often leads to fragmented communication and disproportionate aid distribution \cite{cayamanda2021}. By implementing this centralized mathematical model, local authorities can effectively neutralize the localized Bullwhip Effect. This ensures that scarce resources are systematically allocated based on mathematical equity, guaranteeing that life-sustaining aid reaches the most isolated and severely affected communities in Davao City, rather than being redundantly concentrated in highly accessible urban centers.

\section{Scope and Limitations}
This investigation is focused strictly on the assignment and allocation of resources utilizing an exact mathematical approach. It does not address the physical routing, delivery mechanics, or transportation logistics of those resources. The system's architecture is limited to a Hub-and-Spoke centralized decision framework designed to prevent the Bullwhip Effect computationally. The core engine utilizes a deterministic Mixed-Integer Linear Programming (MILP) model. While classical stochastic programming uses 'two-stage' to denote recourse actions under uncertainty, in this study the two-stage structure refers strictly to a deterministic, two-tier hierarchical decision architecture: Stage 1 establishes tactical facility activation ($y_i \in \{0, 1\}$) under administrative capacity limits, and Stage 2 computes operational continuous resource dispatch ($x_{ijr} \ge 0$), solved simultaneously within a unified optimization formulation

Furthermore, while resource allocation can be evaluated through various cardinal and ordinal criteria, this model explicitly distinguishes between its optimization objective and its evaluation metrics. The model utilizes only one active optimization objective: the minimization of population-weighted deprivation costs. All other fairness metrics detailed in this study—specifically Max Utilitarian Efficiency, Max-Min Fair Share, Max Leximin, Max Nash, and Envy-Freeness—are not enforced as simultaneous mathematical constraints during the solve. Instead, they serve strictly as post-hoc evaluation metrics to analytically assess the equity profile of the optimal allocation matrix generated by the primary objective.

\section{Definition of Terms}
\begin{description}
\item[Bullwhip Effect] – a phenomenon in supply chain management where a lack of communication and coordination leads to the overestimation or underestimation of needs, resulting in either wasted resources or critical shortages.
\item[Cardinal Allocation Criteria] – fairness evaluation metrics that rely on absolute, quantifiable values of utility, including max utilitarian, pure egalitarian, max leximin, and Max Nash.
\item[Deprivation Cost Functions] – mathematical models utilized to measure and represent the increase in human suffering over time, serving as a foundational proxy for evaluating the equity of allocation outcomes.
\item[Disaster Risk Reduction and Management Operations Center (DRRMOC)] – a localized, centralized coordinating body responsible for directing, evaluating, and managing disaster response and relief operations within a specific local government unit or municipality.
\item[Equity] – an ethical framework for allocation that prioritizes the assignment of aid based on the vulnerability of the population and the severity of the disaster's impact, as opposed to providing uniform, equal amounts to everyone.
\item[Hub-and-Spoke Model] – a centralized decision-making architecture where a central authority (the hub) aggregates data and determines the mathematical allocation of goods to various relief points (the spokes) rather than managing the physical routing.
\item[Max Leximin] – a cardinal allocation criterion that focuses on maximizing the utility of the worst-off areas; if the worst-off are equal, it maximizes the utility of the second worst-off, and so forth.
\item[Max-Min Fairness] – an ordinal allocation strategy that seeks to maximize the utility or relief assigned to the most disadvantaged or worst-off areas first.
\item[Max Utilitarian] – a cardinal allocation criterion that seeks to maximize the total overall utility or benefit across all areas, regardless of how unequally it might be divided.
\item[Mixed-Integer Linear Programming (MILP)] – an exact mathematical optimization technique used to determine the best possible outcome by maximizing or minimizing a linear objective function, subject to a set of strict linear constraints where specific variables are restricted to integer values.
\item[Multi-Criteria Fairness] – A comprehensive evaluation approach that assesses resource allocations against multiple ethical dimensions (such as Max-Min fairness, Nash equity, and utilitarian efficiency) post-optimization, while relying on a population-weighted deprivation cost objective to drive the primary allocation.
\item[Ordinal Allocation Criteria] – fairness evaluation metrics that assess allocations based on relative preferences and rankings, including Pareto-efficiency, envy-freeness, and max-min fair share.
\item[Pure Egalitarian] – a cardinal allocation criterion that strictly requires all entities or areas to receive the exact same level of utility or resources, focusing purely on strict equality over varying needs.
\item[Resource Allocation Problem (RAP)] – the core mathematical and logical challenge during a crisis, defined as the assignment of scarce life-sustaining resources under intense time pressure and environmental uncertainty.
\item [Two-Stage Deterministic Hierarchical Structure] – a mathematical decision structure that organizes decision variables into two logical levels—tactical hub activation ($y_i$) in the first tier and operational commodity allocation ($x_{ijr}$) in the second tier—which are solved simultaneously in a single deterministic program, distinct from multi-period stochastic recourse formulations.
\end{description}